\subsection{是同一理论}

让 \(\Sigma\) 并且 \(\Theta\) 中的两个结构物种，具有相同的主基集 \(\mathfrak{C}\) 分别是物种 \(x_{1},\ldots ,x_{n}\) . 设 \(s,t\) 的泛结构。假设满足以下条件： \(\Sigma,\Theta\) 关系

\begin{enumerate}
\def\labelenumi{(\arabic{enumi})}
\item
  我们有一个演绎过程 \(\mathbf{P}\{x_1, \ldots, x_n, s\}\) ，把物种 \(\Theta\) 在 \(x_1, \ldots, x_n\) 上的结构由物种 \(\Sigma\) 在 \(x_1, \ldots, x_n\) 上的结构演绎出来。
\item
  我们有一个演绎过程 \(\mathbb{Q}\{x_1, \ldots, x_n, t\}\) ，把物种 \(\Sigma\) 在 \(x_1, \ldots, x_n\) 上的结构由物种 \(\Theta\) 在 \(x_1, \ldots, x_n\) 上的结构演绎出来。
\item
  ，以及关系 \(\mathbf{Q}\{x_1, \ldots, x_n, \mathbf{P}\{x_1, \ldots, x_n, s\}\} = s\) 是定理于 \(\mathfrak{G}_{\Sigma}\) ，关系B \(\mathbf{P}\{x_1, \ldots, x_n, \mathbf{Q}\{x_1, \ldots, x_n, t\}\} = t\) 是定理于 \(\mathfrak{G}_{\Theta}\) .
\end{enumerate}

结构的物种 \(\Sigma\) 并且 \(\Theta\) 则称通过演绎程序P和Q是等价的。在这种情况下，对于每一个定理B \(\{x_1, \ldots, x_n, s\}\) 在该理论中 \(\mathfrak{G}_{\Sigma}\) ；反之，对于每个定理C \(\{x_1, \ldots, x_n, Q\}\) 是定理于 \(\mathfrak{G}_{\Theta}\) ，关系C \(\{x_1, \ldots, x_n, t\}\) 在该理论中 \(\mathfrak{G}_{\Theta}\) 如果U是物种 \(\{x_1, \ldots, x_n, P\}\) 是定理于 \(\mathfrak{G}_{\Sigma}\) .

的一个结构，则由程序P从U演绎出的结构称为与U等价。准则CST6蕴含以下结论： \(\Sigma\) CST7. 设

分别是物种 \(\mathcal{S},\mathcal{S}'\) 的两个结构。设 \(\Sigma\) 在主基集上 \((\mathbf{E}_1,\ldots ,\mathbf{E}_n),(\mathbf{E}_1',\ldots ,\mathbf{E}_n')\) 是物种 \(\mathcal{S}_0,\mathcal{S}_0'\) 的结构，它们分别与 \(\Theta\) 等价。 \(\mathcal{S}\) 并且 \(\mathcal{S}'\) 。为了使得 \((g_1,\dots,g_n)\) 是关于结构 \(\mathcal{S}_0\) 并且 \(\mathcal{S}_0'\) 的一个同构，当且仅当 \((g_1,\dots,g_n)\) 是关于结构 \(\mathcal{S}\) 并且 \(\mathcal{S}'\) .

在实践中，我们不区分两个等价结构物种的理论 \(\mathfrak{G}_{\Sigma}\) 并且 \(\mathfrak{G}_{\Theta}\) 。

\minorheading{示例}

*（1）设 \(\Sigma\) 是交换群结构的物种； \(\Sigma\) 有一个（主）基集A，其泛结构由单个字母F组成； \(\Sigma\) 是 \(F \in \mathfrak{P}((A \times A) \times A)\) 的典型刻画，并且我们将 \(\Sigma\) 的公理记为R \{A, F\}。该公理特别蕴含F是某个函数的图（群的“合成律”；参见第4节，例2）。在理论 \(\mathfrak{G}_{\Sigma}\) （其中 \(\mathfrak{G}\) 表示集合论）中，我们定义一个项M\{A, F\} ，它是 \(\mathfrak{P}((Z \times A) \times A)\) 中的一个函数图，并满足以下关系B\{M, A, F\}:

\((\forall x)(\forall y)(\forall n)((x \in A \text{ 且 } y \in A \text{ 且 } n \in \mathbf{Z})\)

\[
\Longrightarrow (\mathbf {M} (n, \mathrm{F} (x, y)) = \mathrm{F} (\mathbf {M} (n, x), \mathrm{M} (n, y))))
\]

并且 \((\forall x)(\forall m)(\forall n)((x\in A\) 并且 \(m\in \mathbf{Z}\) 并且 \(n\in \mathbf{Z})\)

\[
\Rightarrow (\mathrm{M} (m + n, x) = \mathrm{F} (\mathrm{M} (m, x), \mathrm{M} (n, x))))
\]

并且 \((\forall x)(\forall m)(\forall n)((x\in A\) 并且 \(m\in \mathbf{Z}\) 并且 \(n\in \mathbf{Z})\)

\[
\Longrightarrow (\mathbf {M} (m, \mathbf {M} (n, x)) = \mathbf {M} (m n, x)))
\]

并且

\[
(\forall x) ((x \in A) \Longrightarrow (M (1, x) = x)).
\]

（“A 的元素与整数相乘”）。

考虑该物种 \(\Theta\) 的 \(\mathbf{Z}\) -模结构，其具有单一的主基集 A，其中 \(\mathbf{Z}\) 作为辅助集合，其一般结构包含两个字母G、L，具有典型的刻画

\[
\mathbf {G} \in \mathfrak {P} ((\mathbf {A} \times \mathbf {A}) \times \mathbf {A}) \quad \text { 且 } \quad \mathbf {L} \in \mathfrak {P} ((\mathbf {Z} \times \mathbf {A}) \times \mathbf {A})
\]

以及公理

\[
\text{``}R\{A,G\}\text{ 且 }(L\text{ 是函数图})\text{ 且 }B\{L,A,G\}\text{''.}
\]

立即验证可知，项F、M构成了一种物种结构的演绎程序。 \(\Theta\) 到物种结构的演绎程序 \(\Sigma\) ，并且项 G 是物种结构的一个演绎过程 \(\Sigma\) 到物种结构的演绎程序 \(\Theta\) 。此外，上述条件（3）平凡地成立。因此我们可以说，交换群结构的物种与 Z-模结构的物种是等价的。

\begin{enumerate}
\def\labelenumi{(\arabic{enumi})}
\setcounter{enumi}{1}
\tightlist
\item
  让 \(\Sigma\) 是拓扑结构的种类（第4号，例3），设A为（主）基集，并设V为 \(\Sigma\) 。考虑关系
\end{enumerate}

\[
x \in A \text {   且   } X \subset A \text {   且   } (\forall U) ((U \in V \text {   且   } x \in U) \Longrightarrow (X \cap U \neq \emptyset)).
\]

这个关系有一个图。 \(P \subset \mathfrak{P}(A) \times A\) 关于这一对 \((X, x)\) ; \(P\{A, V\}\) 是一个称为“所有二元组的集合”的术语 \((X, x)\) 使得x位于X关于拓扑V''的闭包中。然后我们可以证明（参见一般拓扑学，第一章，§1）以下关系是定理 \(G_{\Sigma}\) :

\[
\begin{array}{c} \mathrm{P} (\phi) = \phi , \\ (\forall \mathrm{Y}) ((\mathrm{Y} \subset \mathrm{A}) \Longrightarrow (\mathrm{Y} \subset \mathrm{P} (\mathrm{Y}))), \\ (\forall \mathrm{Y}) ((\mathrm{Y} \subset \mathrm{A}) \Longrightarrow (\mathrm{P} (\mathrm{P} (\mathrm{Y})) = \mathrm{P} (\mathrm{Y}))), \\ (\forall \mathrm{Y}) (\forall \mathrm{Z}) ((\mathrm{Y} \subset \mathrm{A} \text {且} \mathrm{Z} \subset \mathrm{A}) \Longrightarrow (\mathrm{P} (\mathrm{Y} \cup \mathrm{Z}) = \mathrm{P} (\mathrm{Y}) \cup \mathrm{P} (\mathrm{Z}))). \end{array}
\]

考虑结构的物种。 \(\Theta\) ，具有一个单一（主）基集 A，其一般结构由单个字母 W 组成，其典型特征为 \(W \in \mathfrak{P}(\mathfrak{P}(A) \times A)\) 并且作为公理

\[
\begin{array}{c} \mathrm{W} (\emptyset) = \emptyset \quad \text {且} \quad (\forall \mathrm{Y}) ((\mathrm{Y} \subset \mathrm{A}) \Longrightarrow (\mathrm{Y} \subset \mathrm{W} (\mathrm{Y}))) \\ \text {且} \quad (\forall \mathrm{Y}) ((\mathrm{Y} \subset \mathrm{A}) \Longrightarrow (\mathrm{W} (\mathrm{W} (\mathrm{Y})) = \mathrm{W} (\mathrm{Y}))) \\ \text {且} \quad (\forall \mathrm{Y}) (\forall \mathrm{Z}) ((\mathrm{Y} \subset \mathrm{A} \quad \text {且} \quad \mathrm{Z} \subset \mathrm{A}) \Longrightarrow (\mathrm{W} (\mathrm{Y} \cup \mathrm{Z}) = \mathrm{W} (\mathrm{Y}) \cup \mathrm{W} (\mathrm{Z}))). \end{array}
\]

同时也要考虑这个关系

\[
\mathrm{U} \subset \mathrm{A} \quad \text { 且 } \quad (\forall x) ((x \in \mathrm{U}) \Longrightarrow x \notin \mathrm{W} (\mathrm{A} - \mathrm{U})).
\]

所有满足此关系的集合是 \(U \in \mathfrak{P}(A)\) 一个子集 \(Q\{A, W\}\) 的 \(\mathfrak{P}(A)\) . 然后我们可以证明（一般拓扑学，第一章，第1节，习题10）以下关系是定理 \(G_{\Theta}\) :

\[
\begin{array}{c} \mathrm{A} \in \mathrm{Q}, \\ (\forall \mathrm{M}) ((\mathrm{M} \subset \mathrm{Q}) \Longrightarrow \left(\left(\bigcup_ {\mathrm{X} \in \mathrm{M}} \mathrm{X}\right) \in \mathrm{Q}\right), \end{array}
\]

\[
(\forall \mathrm{X}) (\forall \mathrm{Y}) ((\mathrm{X} \in \mathrm{Q} \text {   且   } \mathrm{Y} \in \mathrm{Q}) \Longrightarrow ((\mathrm{X} \cap \mathrm{Y}) \in \mathrm{Q})).
\]

因此，这些项 \(P\{A,V\}\) 并且 \(Q\{A,W\}\) 满足上述条件（1）和（2），并且容易看出它们也满足条件（3）。结构种类 \(\Sigma\) 并且 \(\Theta\) 因此是等价的，因此我们将每一种结构种类 \(\Theta\) 视为一种拓扑，即根据演绎过程与之对应的那种拓扑 \(Q\{A,W\}\). *

